\begin{abstract}
Open-set text recognition in Chinese historical texts remains challenging because of a vast array of ancient characters and a pronounced long-tail distribution, which leaves minimal samples for many rare tail classes. Mitigating this imbalance crucially helps recognize novel characters often hidden within these underrepresented categories. Moreover, Chinese characters share structural parts, making localized feature enhancement highly beneficial for both head and tail classes alike. We therefore propose SpindleNet, which bolsters local component representation by embedding a Radical Residual Module and expanding the channel capacity of middle layers to capture fine-grained details. To constrain parameter growth and complexity, we shorten the depth of subsequent layers, resulting in a distinctive "spindle" shape. Extensive evaluations on three difficult Chinese ancient book datasets, TKH, MTH1000, and MTH1200, show that SpindleNet significantly surpasses existing methods, offering improved feature extraction for tail-class characters and thereby establishing a new state-of-the-art performance across varied tasks.
\end{abstract}
